\documentclass[a4paper,10pt]{article}
\usepackage{lamia-pedagogy}

\begin{document}

% =========================================================================
% PAGE 1 : Fiche de Cours - Définition & Synonymie
% =========================================================================
\ficheHeader{1 Heure}{\makebox[2.5cm]{\dotfill}}{1\textsuperscript{ère} AM}{2026 / 2027}

\vspace{-3pt}
\competenceBanner{\faBookOpen\ Projet 01 / Séquence 01 — Séance 04 : Vocabulaire}

\vspace{-2pt}
\modaliteSupportBar{Observation collective puis entraînement individuel}{Support : Le brossage des dents / Cahier de cours}

\vspace{1pt}

% Texte d'observation
\begin{tcolorbox}[
  enhanced,
  colback=primarylight!40,
  colframe=primary,
  arc=2mm,boxrule=0.8pt,
  top=4pt,bottom=4pt,left=8pt,right=8pt,
  title={\small\textbf{\textcolor{white}{\faEye\ Texte support d'observation}}},
  coltitle=white,colbacktitle=primary
]
{\small
\guillemotleft\ Le brossage des dents \textbf{est} une pratique qui assure une bonne hygiène bucco-dentaire. Il \textbf{consiste à} utiliser une brosse à dents et du dentifrice pour nettoyer les dents, \textbf{éliminer} les résidus d'aliments et préserver l'émail dentaire. \guillemotright
}
\end{tcolorbox}

\vspace{3pt}

% Règles mémos
\begin{minipage}[t]{0.485\textwidth}
\begin{tcolorbox}[
  enhanced,
  colback=secondarylight!50,
  colframe=secondary,
  arc=2mm,boxrule=0.8pt,
  top=4pt,bottom=4pt,left=7pt,right=7pt,
  title={\small\textbf{\textcolor{secondary!90!black}{\faStar\ 1. La Définition}}},
  colbacktitle=white,coltitle=secondary!90!black,
  attach boxed title to top left={yshift=-2mm,xshift=4mm},
  boxed title style={boxrule=0.6pt,colframe=secondary,arc=1.2mm,top=1pt,bottom=1pt}
]
\vspace{2pt}
{\small
\textbf{Qu'est-ce qu'une définition ?}\\
C'est un procédé qui donne le sens exact d'un mot ou d'un concept.

\vspace{3pt}
\textbf{Verbes introducteurs courants :}
\begin{itemize}[leftmargin=*,noitemsep,topsep=1pt]
  \item \textbf{être} (\textit{Le savon \textbf{est} un produit...})
  \item \textbf{consister à} (\textit{Il \textbf{consiste à} frotter...})
  \item \textbf{désigner / signifier}
\end{itemize}
}
\end{tcolorbox}
\end{minipage}%
\hfill
\begin{minipage}[t]{0.485\textwidth}
\begin{tcolorbox}[
  enhanced,
  colback=accentlight!35,
  colframe=accent,
  arc=2mm,boxrule=0.8pt,
  top=4pt,bottom=4pt,left=7pt,right=7pt,
  title={\small\textbf{\textcolor{accent!90!black}{\faStar\ 2. La Synonymie}}},
  colbacktitle=white,coltitle=accent!90!black,
  attach boxed title to top left={yshift=-2mm,xshift=4mm},
  boxed title style={boxrule=0.6pt,colframe=accent,arc=1.2mm,top=1pt,bottom=1pt}
]
\vspace{2pt}
{\small
\textbf{Qu'est-ce qu'un synonyme ?}\\
C'est un mot de \textbf{même sens} (ou sens très voisin) et de \textbf{même nature grammaticale}.

\vspace{3pt}
\textbf{Exemples clés :}
\begin{itemize}[leftmargin=*,noitemsep,topsep=1pt]
  \item \textbf{Verbes :} éliminer = supprimer = bannir
  \item \textbf{Verbes :} nettoyer = laver
  \item \textbf{Noms :} une pratique = une habitude
  \item \textbf{Adjectifs :} indispensable = essentiel
\end{itemize}
}
\end{tcolorbox}
\end{minipage}

\vspace{3pt}

% Schéma conceptuel
\begin{tcolorbox}[
  enhanced,
  colback=white,
  colframe=cardborder,
  arc=1.5mm,boxrule=0.6pt,
  top=4pt,bottom=4pt,left=7pt,right=7pt,
  title={\small\textbf{\textcolor{primary}{\faProjectDiagram\ Schéma de synthèse}}},
  coltitle=primary,colbacktitle=primarylight,titlerule=0.4pt
]
\centering
\begin{tikzpicture}[scale=0.9, every node/.style={transform shape}]
  \node[draw=secondary, rounded corners, fill=secondarylight, thick, text=primary, align=center, inner sep=5pt] (mot) at (0,0) {\textbf{Mot cible}\\\textit{(ex: Éliminer)}};
  \node[draw=accent, rounded corners, fill=accentlight!40, thick, text=accent!90!black, align=center, inner sep=5pt] (syn) at (5,0) {\textbf{Synonyme}\\\textit{(ex: Supprimer)}};
  \draw[<->, thick, primary] (mot) -- node[above]{\small Même sens} node[below]{\small Même classe grammaticale} (syn);
\end{tikzpicture}
\end{tcolorbox}

\newpage

% =========================================================================
% PAGE 2 : Fiche d'Exercices d'Application
% =========================================================================
\ficheHeader{1 Heure}{Nom : \makebox[4cm]{\dotfill}}{1\textsuperscript{ère} AM}{Classe : \makebox[1.5cm]{\dotfill}}

\vspace{-3pt}
\competenceBanner{\faPenFancy\ Fiche d'Exercices : Définition et Synonymie}

\vspace{2pt}

% ACTIVITÉ 1 : Repérage de définition
\begin{tcolorbox}[
  enhanced,
  colback=white,
  colframe=cardborder,
  arc=2mm,boxrule=0.8pt,
  top=4pt,bottom=4pt,left=8pt,right=8pt,
  title={\small\textbf{\textcolor{primary}{\faHighlighter\ Activité 01 : Souligne la définition et entoure le verbe introducteur}} \hfill \badge[secondary]{La Définition}},
  coltitle=primary,colbacktitle=primarylight,titlerule=0.5pt,titlerule style=cardborder
]
{\small
\begin{enumerate}[label=\textbf{\arabic*.},leftmargin=*,noitemsep,topsep=2pt]
  \item La carie dentaire \textbf{est} une infection des tissus de la dent provoquée par les bactéries.
  \item Les ablutions \textbf{consistent à} se laver soigneusement le visage, les mains et les pieds avant la prière.
  \item Le savon liquide \textbf{désigne} un produit nettoyant utilisé pour l'hygiène quotidienne.
\end{enumerate}
}
\end{tcolorbox}

\vspace{2pt}

% ACTIVITÉ 2 : Tableau d'association des synonymes
\begin{tcolorbox}[
  enhanced,
  colback=white,
  colframe=cardborder,
  arc=2mm,boxrule=0.8pt,
  top=4pt,bottom=4pt,left=8pt,right=8pt,
  title={\small\textbf{\textcolor{primary}{\faLink\ Activité 02 : Relie chaque mot de la colonne A à son synonyme dans la colonne B}} \hfill \badge[accent]{Synonymie}},
  coltitle=primary,colbacktitle=primarylight,titlerule=0.5pt,titlerule style=cardborder
]
{\small
\begin{center}
\begin{tabularx}{0.88\linewidth}{|X|c|X|}
\hline
\rowcolor{primarylight}\textbf{Colonne A} & $\longleftrightarrow$ & \textbf{Colonne B (Synonyme)} \\
\hline
1. Éliminer & $\bullet\quad\bullet$ & a. Indispensable \\
\hline
2. Nettoyer & $\bullet\quad\bullet$ & b. Supprimer \\
\hline
3. Essentiel & $\bullet\quad\bullet$ & c. Habitude \\
\hline
4. Pratique & $\bullet\quad\bullet$ & d. Laver \\
\hline
\end{tabularx}
\end{center}
}
\end{tcolorbox}

\vspace{2pt}

% ACTIVITÉ 3 : Réécriture de phrases
\begin{tcolorbox}[
  enhanced,
  colback=white,
  colframe=cardborder,
  arc=2mm,boxrule=0.8pt,
  top=4pt,bottom=4pt,left=8pt,right=8pt,
  title={\small\textbf{\textcolor{primary}{\faSync*\ Activité 03 : Remplace le mot en gras par son synonyme}} \hfill \badge[emerald]{Réécriture}},
  coltitle=primary,colbacktitle=primarylight,titlerule=0.5pt,titlerule style=cardborder
]
{\small
\begin{enumerate}[label=\textbf{\arabic*.},leftmargin=*,noitemsep,topsep=2pt]
  \item Le savon permet d'\textbf{éliminer} les microbes. (\textit{supprimer / salir})\\
    $\longrightarrow$ Le savon permet de \drule{6cm} les microbes.
  \item Le brossage des dents est un geste \textbf{indispensable}. (\textit{inutile / nécessaire})\\
    $\longrightarrow$ Le brossage des dents est un geste \drule{6cm}.
\end{enumerate}
}
\end{tcolorbox}

\vspace{2pt}

% ACTIVITÉ 4 : Production d'une définition
\begin{tcolorbox}[
  enhanced,
  colback=white,
  colframe=cardborder,
  arc=2mm,boxrule=0.8pt,
  top=4pt,bottom=4pt,left=8pt,right=8pt,
  title={\small\textbf{\textcolor{primary}{\faPenSquare\ Activité 04 : Rédige une définition}} \hfill \badge[softred]{Production}},
  coltitle=primary,colbacktitle=primarylight,titlerule=0.5pt,titlerule style=cardborder
]
{\small
À partir des mots suivants, construis une phrase contenant une définition :\\
\textbf{Mots :} \textit{Le dentifrice / être / une pâte nettoyante / protéger les dents contre les caries.}\\
$\longrightarrow$ \drule{\linewidth}
}
\end{tcolorbox}

\end{document}
